\clearpage
\sectionkicker{Source-backed technical notes}
\subsection*{Multimodal — 生成からリアルタイム統合へ: Technical Notes}
この欄は記事本文で圧縮した一次資料上の情報を、比較・再検証しやすい形へ展開したものである。新しい外部情報は追加せず、Selection済みEvidenceのchronology、normalized claim、limitations、source URLのみを再配置する。
\medskip
\subsection*{Theme at a glance}
\addcontentsline{toc}{subsection}{Theme at a glance}
\begin{center}
\footnotesize
\begin{tabularx}{\linewidth}{@{}p{0.20\linewidth}p{0.10\linewidth}p{0.14\linewidth}X@{}}
\toprule
Artifact & Role & Type & Objective chronology \\
\midrule
MiniMax H3 & PRIMARY & MODEL & 2026-07-31 (MODEL\_RELEASE) \\
GPT-Live & PRIMARY & MODEL & 2026-07-08 (MODEL\_RELEASE) \\
Alibaba Model Studio July multimodal updates & PRIMARY & MODEL UPDATE & 2026-07-01 (MODEL\_UPDATE); 2026-07-14 (MODEL\_UPDATE) \\
\bottomrule
\end{tabularx}
\end{center}
\normalsize

\begin{technicalnote}{MiniMax H3}{PRIMARY}
\begin{tabularx}{\linewidth}{@{}>{\bfseries}p{0.22\linewidth}X@{}}
Organization & MiniMax \\
Artifact type & MODEL \\
Chronology & 2026-07-31 (MODEL\_RELEASE) \\
\end{tabularx}
\smallskip
{\bfseries 一次資料から整理したtechnical points}
\begin{itemize}[leftmargin=1.5em,itemsep=0.35em]
\item \textbf{一次情報で確認できる事実}: MiniMax announced H3 at the end of July as a multimodal generative model intended to unify tasks across image, video, and audio generation and cross-modal context understanding.
\item \textbf{Vendor claim}: MiniMax describes H3 as supporting up to 2K video output with native stereo sound and multimodal references spanning image, audio, and video.
\item \textbf{一次情報で確認できる事実}: At announcement time MiniMax said it planned to open the model weights in the coming days, so the launch itself must not be written as an already-completed open-weight release.
\end{itemize}
{\bfseries 読む際の境界}
\begin{itemize}[leftmargin=1.5em,itemsep=0.35em]
\item \textbf{分析上の留意点}: Generation quality and capability examples are vendor demonstrations; the announcement explicitly described weight publication as planned rather than complete at that moment.
\end{itemize}
{\bfseries Primary source}
\begingroup\sloppy
\begin{itemize}[leftmargin=1.5em,itemsep=0.25em]
\item {\scriptsize\url{https://www.minimax.io/news}}
\end{itemize}
\endgroup
{\scriptsize\color{SurveyMuted}Source-bound record: \texttt{evidence:SP-2026-M07:item-official-index-minimax-news-8be7d33438}.}
\end{technicalnote}
\medskip

\begin{technicalnote}{GPT-Live}{PRIMARY}
\begin{tabularx}{\linewidth}{@{}>{\bfseries}p{0.22\linewidth}X@{}}
Organization & OpenAI \\
Artifact type & MODEL \\
Chronology & 2026-07-08 (MODEL\_RELEASE) \\
\end{tabularx}
\smallskip
{\bfseries 一次資料から整理したtechnical points}
\begin{itemize}[leftmargin=1.5em,itemsep=0.35em]
\item \textbf{一次情報で確認できる事実}: OpenAI announced GPT-Live on July 8, 2026 as a real-time model designed for simultaneous listening and speaking.
\item \textbf{Vendor claim}: OpenAI describes GPT-Live as delegating deeper work to a frontier model behind the live interaction loop, with GPT-5.5 used at launch.
\end{itemize}
{\bfseries 読む際の境界}
\begin{itemize}[leftmargin=1.5em,itemsep=0.35em]
\item \textbf{分析上の留意点}: Availability and latency/capability statements originate from OpenAI; independent real-time quality evaluation is not established by this source.
\end{itemize}
{\bfseries Primary source}
\begingroup\sloppy
\begin{itemize}[leftmargin=1.5em,itemsep=0.25em]
\item {\scriptsize\url{https://openai.com/index/introducing-gpt-live}}
\end{itemize}
\endgroup
{\scriptsize\color{SurveyMuted}Source-bound record: \texttt{evidence:SP-2026-M07:item-official-feed-3fb14968493dc68260e4-dd57e9af7c}.}
\end{technicalnote}
\medskip

\begin{technicalnote}{Alibaba Model Studio July multimodal updates}{PRIMARY}
\begin{tabularx}{\linewidth}{@{}>{\bfseries}p{0.22\linewidth}X@{}}
Organization & Alibaba Cloud \\
Artifact type & MODEL UPDATE \\
Chronology & 2026-07-01 (MODEL\_UPDATE); 2026-07-14 (MODEL\_UPDATE) \\
\end{tabularx}
\smallskip
{\bfseries 一次資料から整理したtechnical points}
\begin{itemize}[leftmargin=1.5em,itemsep=0.35em]
\item \textbf{一次情報で確認できる事実}: Alibaba Cloud Model Studio lists international availability of Qwen-Audio 3.0 TTS models on July 14, 2026.
\item \textbf{一次情報で確認できる事実}: The same lifecycle page lists a Wan 2.7 reference-to-video snapshot for the international region on July 1, 2026.
\item \textbf{Vendor claim}: Alibaba describes the Qwen-Audio 3.0 Flash variant as targeting low-latency real-time interaction and the Plus variant as targeting higher-quality professional scenarios.
\end{itemize}
{\bfseries 読む際の境界}
\begin{itemize}[leftmargin=1.5em,itemsep=0.35em]
\item \textbf{分析上の留意点}: Availability is region-specific and the performance/quality descriptions are vendor claims from Alibaba Cloud Model Studio.
\end{itemize}
{\bfseries Primary source}
\begingroup\sloppy
\begin{itemize}[leftmargin=1.5em,itemsep=0.25em]
\item {\scriptsize\url{https://www.alibabacloud.com/help/en/model-studio/newly-released-models}}
\end{itemize}
\endgroup
{\scriptsize\color{SurveyMuted}Source-bound record: \texttt{evidence:SP-2026-M07:item-official-index-alibaba-model-studio-154c528d30}.}
\end{technicalnote}
\medskip
